
\clearpage
\chapter{ארכיטקטורת \textenglish{xLSTM}}

\section{מוטיבציה ורקע}

\textenglish{xLSTM} \cite{beck2024xlstm} נולד מתוך השאלה: האם ניתן לשפר \textenglish{LSTM} קלאסי \cite{hochreiter1997lstm} עד כדי תחרות עם \textenglish{Transformer} מבלי להוסיף מורכבות ריבועית? שני חידושים עיקריים מאפשרים זאת:

\subsection{\textenglish{sLSTM}: שערי אקספוננציאל}

במקום שער סיגמואיד קלאסי, \textenglish{sLSTM} משתמש בפונקציית שכחה אקספוננציאלית:

\begin{equation}
f_t = \exp(-\exp(w_f h_{t-1} + r_f x_t + b_f))
\end{equation}

גישה זו מאפשרת "שכחה מהירה" אדפטיבית ופותרת בעיות כמו \textenglish{over-counting} בזיכרון ארוך-טווח.

\subsection{\textenglish{mLSTM}: זיכרון מטריצה}

\textenglish{mLSTM} מרחיב את תא ה-\textenglish{LSTM} לייצוג מטריצה:

\begin{equation}
C_t = f_t \odot C_{t-1} + i_t \cdot v_t k_t^T \in \mathbb{R}^{d \times d}
\end{equation}

כאשר \(v_t\), \(k_t\) הם וקטורי ערך ומפתח. מטריצה זו מאפשרת אחסון מגוון רב יותר של מידע בהשוואה לוקטור הסמוי הקלאסי.

\begin{figure}[H]
\centering
\includegraphics[width=0.85\textwidth]{assets/architecture.png}
\caption{\textenglish{Architecture Comparison: PatchTST vs xLSTM (Python-generated)}}
\label{fig:arch}
\end{figure}
